% Unit 068 English translation of chapter5.tex lines 578--739.
% Source: Methods of Algebra, Volume 2, commit
% 9a5803ff77dd3257484cb177f851a73770a59dd3 (CC BY 4.0).
% stable-unit-id: o014.aljabr2.chapter5.filtered-complex-spectral-sequence
% Coverage: Complete Section 5.5.

% segment-id: o014.li2.chapter5.filtered-complex-spectral-sequence.h001
\section{The spectral sequence of a filtered complex}\label{sec:filtered-cplx-ss}
\hypertarget{o014.aljabr2.chapter5.filtered-complex-spectral-sequence}{ }
\index{complex!filtered}
\index{spectral sequence!of a filtered complex}

% segment-id: o014.li2.chapter5.filtered-complex-spectral-sequence.p001
We continue the line of thought from \S\ref{sec:filtered-diff-ss}, still with
a fixed Abelian category $\mathcal{A}$. The result of
Proposition~\ref{prop:filtered-diff-E} for filtered differential objects can
be extended to allow the morphism $d$ to have a degree. In particular, this
theory applies to filtered complexes.

% segment-id: o014.li2.chapter5.filtered-complex-spectral-sequence.p002
Concretely, replace the previous $\mathcal{A}$ by $\mathcal{A}^{\Z}$ and
denote its translation functor by $T$. By Example~\ref{eg:cplx-as-diff}, an
object $(X,d)$ of $(\mathcal{A}^{\Z},T)_d$ can be identified with a complex
$X\in\Obj(\cate{C}(\mathcal{A}))$. If we further take a decreasing filtration
$(X,d,\mathrm{F}^\bullet X)$, then
\begin{align*}
	\Hm(X,d)&=\left(\Hm^n(X)\right)_{n\in\Z}, \\
	\Hm\left(\mathrm{F}^pX,d\right)
	&=\left(\Hm^n\left(\mathrm{F}^pX\right)\right)_{n\in\Z}, \\
	\mathrm{F}^p\Hm^n(X)&:=
	\Image\left[\mathrm{F}^pX^n\cap\Ker(d_X^n)\to\Hm^n(X)\right]
	\quad\text{(Definition~\ref{def:induced-filtration-H}).}
\end{align*}

% segment-id: o014.li2.chapter5.filtered-complex-spectral-sequence.p003
Thus the $E_r$ in the spectral sequence $\mathscr{E}$ constructed above in
fact takes values in $\mathcal{A}^{\Z\times\Z}$ and is a bigraded object.
Besides the filtration degree $p$, there is an ``internal degree'' $n$ coming
from the complex structure; the corresponding translation functors $S$ and
$T$ commute strictly. We already know that $d_r$ has degree $r$ with respect
to $S$; we now determine its degree with respect to $T$.

% segment-id: o014.li2.chapter5.filtered-complex-spectral-sequence.p004
\begin{itemize}
	\item In the exact couple $\mathscr{C}_{(1)}$ used to construct
	$\mathscr{E}$, the morphism $\nwarrow$ has degree $1$ with respect to $T$,
	while the other morphisms have degree zero. Indeed, $\nwarrow$ comes from
	the cohomological connecting morphism and hence has degree $1$ with respect
	to $T$, while the other arrows plainly have degree zero.
	\item Using the description in Lemma~\ref{prop:exact-couple-Z-B}, the same
	statement follows recursively for $\mathscr{C}_{(r)}$ for every $r$.
	\item Consequently, both $d_r$ and the $d$ in $(X,d)$ are morphisms of
	degree $1$ with respect to $T$. This is also immediate from the description
	in Proposition~\ref{prop:differential-exact-couple}.
\end{itemize}

% segment-id: o014.li2.chapter5.filtered-complex-spectral-sequence.p005
For applications, the customary convention is to use the indices $p$ and
$q:=n-p$. Thus
\begin{gather*}
	E_r=(E_r^{p,q})_{(p,q)\in\Z^2}
	=\left(Z_r^{p,q}/B_r^{p,q}\right)_{(p,q)\in\Z^2}, \\
	E_0^{p,q}=\left(\gr^pX\right)^{p+q}, \\
	E_1^{p,q}=\Hm^{p+q}\left(\gr^pX,\gr^pd\right), \\
	d_r=\left(d_r^{p,q}:E_r^{p,q}\to E_r^{p+r,q-r+1}\right)_{(p,q)\in\Z^2}
	:E_r\xrightarrow{(r,-r+1)}E_r.
\end{gather*}
In other words, a filtered complex gives a cohomological bigraded spectral
sequence as in Definition~\ref{def:graded-spectral-sequence}. Dually, a chain
complex with an increasing filtration gives a homological bigraded spectral
sequence.

% segment-id: o014.li2.chapter5.filtered-complex-spectral-sequence.e001
\begin{example}\label{eg:canonically-bdd-ss}
	Consider a filtered complex $(X,d,\mathrm{F}^\bullet X)$. If for every
	$n\in\Z$ one has $\mathrm{F}^0X^n=X^n$ and
	$\mathrm{F}^{n+1}X^n=0$, then $E_0$ lies in the first quadrant
	(Definition~\ref{def:quadrant-ss}). This follows immediately from
	$E_0^{p,q}=(\gr^pX)^{p+q}$.
\end{example}

% segment-id: o014.li2.chapter5.filtered-complex-spectral-sequence.t001
\begin{proposition}\label{prop:filtered-cplx-ss}
	For a complex $X\in\Obj(\cate{C}(\mathcal{A}))$ with decreasing filtration
	$\mathrm{F}^\bullet X$, the associated spectral sequence satisfies
	\begin{align*}
		Z_r^{p,q}&=
		\dfrac{\left(\mathrm{F}^pX^{p+q}\cap
		d^{-1}\mathrm{F}^{p+r}X^{p+q+1}\right)+\mathrm{F}^{p+1}X^{p+q}}
		{\mathrm{F}^{p+1}X^{p+q}}, \\
		B_r^{p,q}&=
		\dfrac{\left(\mathrm{F}^pX^{p+q}\cap
		d\mathrm{F}^{p-r+1}X^{p+q-1}\right)+\mathrm{F}^{p+1}X^{p+q}}
		{\mathrm{F}^{p+1}X^{p+q}},
	\end{align*}
	and $d_r^{p,q}:E_r^{p,q}\to E_r^{p+r,q-r+1}$ is induced by
	\[ d^{-1}\mathrm{F}^{p+r}X^{p+q+1}\xrightarrow{d}
	\mathrm{F}^{p+r}X^{p+q+1}\cap\Ker(d). \]
\end{proposition}
\begin{proof}
	In Proposition~\ref{prop:filtered-diff-E}, include the complex degree
	$n=p+q$ and observe that $d$ has degree $1$ with respect to the complex
	translation functor $T$.
\end{proof}

% segment-id: o014.li2.chapter5.filtered-complex-spectral-sequence.p006
Similarly, \eqref{eqn:filtered-diff-infty} has the bigraded version
\begin{equation}\label{eqn:filtered-cplx-infty}
	E_\infty^{p,q}=\dfrac{Z_\infty^{p,q}}{B_\infty^{p,q}}
	=\dfrac{\bigcap_r\left((\mathrm{F}^pX^{p+q}\cap
	d^{-1}\mathrm{F}^{p+r}X^{p+q+1})+\mathrm{F}^{p+1}X^{p+q}\right)}
	{\bigcup_r\left((\mathrm{F}^pX^{p+q}\cap
	d\mathrm{F}^{p-r+1}X^{p+q-1})+\mathrm{F}^{p+1}X^{p,q}\right)},
\end{equation}
provided the displayed $\bigcup$ and $\bigcap$ exist. In this case
$E_\infty=(E_\infty^{p,q})_{(p,q)\in\Z^2}$.

% segment-id: o014.li2.chapter5.filtered-complex-spectral-sequence.p007
Definition--Proposition~\ref{def:diff-obj-convergence} now becomes a bigraded
statement: for all $p,q$, the object $\gr^p\Hm^{p+q}(X)$ can be canonically
realized as a subquotient of $E_\infty^{p,q}$. The convergence properties of
the spectral sequence are correspondingly refined by the grading.

% segment-id: o014.li2.chapter5.filtered-complex-spectral-sequence.d001
\begin{definition}\label{def:filtered-cplx-convergence}
	\index{spectral sequence!weakly convergent, strongly convergent}
	For a filtered complex $(X,d,\mathrm{F}^\bullet X)$ over $\mathcal{A}$,
	construct the associated cohomologically graded spectral sequence
	$(E_r,d_r)_r$. If the $\bigcap_r$ and $\bigcup_r$ in
	\eqref{eqn:filtered-cplx-infty} exist, then
	$\gr^p\Hm^{p+q}(X)$ can be canonically realized as a subquotient of
	$E_\infty^{p,q}$.
	\begin{itemize}
		\item If $\gr^p\Hm^{p+q}(X)=E_\infty^{p,q}$ for every
		$(p,q)\in\Z^2$, then $(E_r,d_r)_r$ is called \emph{weakly convergent}.
		\item If the sequence is weakly convergent and, for every $n\in\Z$,
		the filtration $\mathrm{F}^\bullet\Hm^n(X)$ is exhaustive and complete,
		then $(E_r,d_r)_r$ is called \emph{strongly convergent}.
	\end{itemize}
\end{definition}

% segment-id: o014.li2.chapter5.filtered-complex-spectral-sequence.c001
\begin{convention}\label{con:ss-conv}
	\index[sym1]{Erpq abuts to H@$E_r^{p,q}\Rightarrow\Hm^{p+q}(X)$}
	Strong convergence of a spectral sequence is also denoted by
	$E_r^{p,q}\Rightarrow\Hm^{p+q}(X)$. The subscript $r$ is usually written
	concretely as $1$, $2$, and so on, depending on which page of the spectral
	sequence we wish to describe.

	More generally, suppose we are given a cohomological bigraded spectral
	sequence $\mathscr{E}$, a filtered graded object
	$(H,\mathrm{F}^\bullet H)$, and isomorphisms
	$E_\infty^{p,q}\simeq\gr^pH^{p+q}$, with
	$\mathrm{F}^\bullet H$ exhaustive and complete. We also denote this
	situation by $E_r^{p,q}\Rightarrow H^{p+q}$. Homological bigraded spectral
	sequences are treated similarly; in particular,
	$E^r_{p,q}\Rightarrow H_{p+q}$ entails
	$E^\infty_{p,q}\simeq\gr_pH_{p+q}$.
\end{convention}

% segment-id: o014.li2.chapter5.filtered-complex-spectral-sequence.p008
The following classical convergence theorem is sufficient for initial
applications. For more general convergence conditions, see \cite{Boa99}.

% segment-id: o014.li2.chapter5.filtered-complex-spectral-sequence.t002
\begin{theorem}[Classical convergence theorem]\label{prop:classical-convergence}
	Consider a filtered complex $(X,d,\mathrm{F}^\bullet X)$. Suppose that for
	every $n\in\Z$,
	\begin{compactitem}
		\item the filtration $\mathrm{F}^\bullet X^n$ is exhaustive
		(Definition~\ref{def:filtration-properties}),
		\item there is an $N=N(n)$ such that $\mathrm{F}^NX^n=0$.
	\end{compactitem}
	Then the associated spectral sequence is strongly convergent
	(Definition--Proposition~\ref{def:diff-obj-convergence}).

	If the hypotheses are strengthened so that the filtration on every $X^n$
	is finite (Definition~\ref{def:filtration}), then the induced filtration on
	$\Hm^n(X)$ is also finite. In this case the associated spectral sequence is
	bounded (Definition~\ref{def:bdd-ss}).
\end{theorem}
\begin{proof}
	By the graded version of Lemma~\ref{prop:induced-filtration-H}, the induced
	filtration $\mathrm{F}^\bullet\Hm^n(X)$ is exhaustive and complete. It
	therefore remains to prove weak convergence.

	We need to recall the proof of
	Definition--Proposition~\ref{def:diff-obj-convergence}, which explains how
	to realize $\gr^p\Hm^{p+q}(X)$ as a subquotient of $E_\infty^{p,q}$ for
	each $(p,q)\in\Z^2$. The key point is
	\begin{equation*}\begin{split}
		\bigcap_r&\left(\left(\mathrm{F}^pX^{p+q}\cap
		d^{-1}\mathrm{F}^{p+r}X^{p+q+1}\right)+\mathrm{F}^{p+1}X^{p+q}\right)
		\\
		&\supset(\mathrm{F}^pX^{p+q}\cap\Ker(d))+\mathrm{F}^{p+1}X^{p+q},\\
		\bigcup_r&\left(\left(\mathrm{F}^pX^{p+q}\cap
		d\mathrm{F}^{p-r+1}X^{p+q-1}\right)+\mathrm{F}^{p+1}X^{p+q}\right)
		\\
		&\subset(\mathrm{F}^pX^{p+q}\cap\Image(d))+\mathrm{F}^{p+1}X^{p+q}.
	\end{split}\end{equation*}
	Proving weak convergence is equivalent to strengthening both inclusions to
	equalities. But when $r\gg0$ relative to $p,q$, the expression inside
	$\bigcap_r$ in the first line is simply
	$(\mathrm{F}^pX^{p+q}\cap\Ker(d))+\mathrm{F}^{p+1}X^{p+q}$, so equality
	holds.

	For the second line, the term $+\mathrm{F}^{p+1}X^{p+q}$ can first be
	moved outside $\bigcup_r$. Recall that $d(\mathrm{F}^\bullet X^n)$ is an
	exhaustive filtration of $d(X^n)$. Hence
	\begin{multline*}
		\bigcup_r\left(\mathrm{F}^pX^{p+q}\cap
		d\mathrm{F}^{p-r+1}X^{p+q-1}\right)\\
		\xlongequal{\because\;\text{\eqref{eqn:exhaustion-intersection}}}
		\mathrm{F}^pX^{p+q}\cap\bigcup_r
		d\mathrm{F}^{p-r+1}X^{p+q-1}\\
		=\mathrm{F}^pX^{p+q}\cap d\bigcup_r
		\mathrm{F}^{p-r+1}X^{p+q-1}
		=\mathrm{F}^pX^{p+q}\cap\Image(d).
	\end{multline*}

	Finally, suppose the filtration on every $X^n$ is finite. Then
	$\mathrm{F}^\bullet\Hm^n(X)$ is naturally bounded. To show that the
	spectral sequence is bounded, fix $r,n$ and take $p+q=n$. By the
	description in Proposition~\ref{prop:filtered-cplx-ss}, when $p\gg0$ the
	numerator of $Z_r^{p,q}$ is $0$, while when $p\ll0$ its denominator is
	$X^n$. Thus only finitely many $(p,q)$ satisfy $E_{r+1}^{p,q}\neq0$.
\end{proof}

% segment-id: o014.li2.chapter5.filtered-complex-spectral-sequence.p009
The dual version is clear and will not be repeated.

% segment-id: o014.li2.chapter5.filtered-complex-spectral-sequence.t003
\begin{corollary}[Exact sequence of low-degree terms]
\label{prop:low-degree-ss}
	Let the filtered complex $(X,d,\mathrm{F}^\bullet X)$ satisfy
	\[ \mathrm{F}^0X^n=X^n,\qquad \mathrm{F}^{n+1}X^n=0,
	\qquad n\in\Z, \]
	as in Example~\ref{eg:canonically-bdd-ss}. The associated spectral sequence
	gives the canonical exact sequence
	\[ 0\to E_2^{1,0}\to\Hm^1(X)\to E_2^{0,1}\xrightarrow{d}E_2^{2,0}
	\to\Hm^2(X). \]
	Dually, for a chain complex $X$ with increasing filtration, if
	$\mathrm{F}_{-1}X=0$ and $\mathrm{F}_nX=X$, there is a canonical exact
	sequence
	\[ \Hm_2(X)\to E^2_{2,0}\xrightarrow{d}E^2_{0,1}\to\Hm_1(X)
	\to E^2_{1,0}\to0. \]
\end{corollary}
\begin{proof}
	Substitute $p=2$ into \eqref{eqn:low-exact-coh} to obtain the exact sequence
	\[ 0\to E_\infty^{0,1}\to E_2^{0,1}\xrightarrow{d}E_2^{2,0}
	\to E_\infty^{2,0}\to0. \]
	By the classical convergence theorem~\ref{prop:classical-convergence},
	$E_\infty^{0,1}\simeq\gr^0\Hm^1(X)$,
	$E_\infty^{1,0}\simeq\gr^1\Hm^1(X)$, and
	$E_\infty^{2,0}\simeq\gr^2\Hm^2(X)$. The hypotheses give
	\begin{gather*}
		\gr^2\Hm^2(X)=\mathrm{F}^2\Hm^2(X),\qquad
		\gr^1\Hm^1(X)=\mathrm{F}^1\Hm^1(X),\\
		\gr^0\Hm^1(X)=\Hm^1(X)/\mathrm{F}^1\Hm^1(X)
		=\Hm^1(X)/E_\infty^{1,0}.
	\end{gather*}
	Moreover, substituting $p=1$ into \eqref{eqn:edge-ss-coh} gives
	$E_2^{1,0}=E_\infty^{1,0}$. Splicing these equalities together yields the
	asserted exact sequence. The homological version is analogous.
\end{proof}

% segment-id: o014.li2.chapter5.filtered-complex-spectral-sequence.p010
If $E_r^{p,q}$ is strongly convergent, then once sufficiently many pages
$(E_r,d_r)$ are known, in principle one can read
$(\gr^p\Hm^{p+q}(X))_{p,q\in\Z}$ from $E_\infty$. Passing from
$(\gr^p\Hm^n(X))_p$ to $\Hm^n(X)$, however, amounts to determining a series
of extensions in $\mathcal{A}$, which is generally difficult unless
$\Hm^n(X)$ splits. If we consider only coarser properties, a spectral sequence
can sometimes provide a concise answer. The following example is essentially
an application of the Euler--Poincar\'e principle.

% segment-id: o014.li2.chapter5.filtered-complex-spectral-sequence.l001
\begin{lemma}\label{prop:finite-degeneration}
	Let $\mathscr{E}=(E_r,d_r)_{r\geq1}$ be a cohomological bigraded spectral
	sequence, and suppose there is an $r$ such that
	\[ \left\{(p,q)\in\Z^2:E_r^{p,q}\neq0\right\}
	\quad\text{is a finite set}. \]
	Then for $r\gg0$, $\mathscr{E}$ degenerates at the $E_r$ page.
\end{lemma}
\begin{proof}
	By the hypothesis, for sufficiently large $r$ all nonzero terms of $E_r$
	are concentrated in a region of finite length and width. Apply
	Example~\ref{eg:SS-finite-width}.
\end{proof}

% segment-id: o014.li2.chapter5.filtered-complex-spectral-sequence.t004
\begin{proposition}\label{prop:EP-ss}
	Let the cohomological bigraded spectral sequence $(E_r,d_r)_r$ satisfy the
	following conditions:
	\begin{compactitem}
		\item there is an $r$ such that
		$\{(p,q)\in\Z^2:E_r^{p,q}\neq0\}$ is a finite set;
		\item there is strong convergence $E_r^{p,q}\Rightarrow H^{p+q}$ in
		the sense of Convention~\ref{con:ss-conv}.
	\end{compactitem}
	Then in the group $\mathrm{K}_0(\mathcal{A})$ of
	Definition~\ref{def:K0}, for $r\gg0$ one has
	\[ \sum_{n\in\Z}(-1)^n[H^n]
	=\sum_{p,q\in\Z}(-1)^{p+q}[E_r^{p,q}], \]
	and both sides are finite sums.
\end{proposition}
\begin{proof}
	For $r\gg0$, only finitely many terms on the right are nonzero. Moreover,
	\[ E_{r+1}^{p,q}\simeq\Hm\left[
	E_r^{p-r,q+r-1}\xrightarrow{d}E_r^{p,q}
	\xrightarrow{d}E_r^{p+r,q-r+1}\right]. \]
	If degree is counted by $n:=p+q$, then $d$ has degree $1$. Take the
	alternating sum over $n$ in $\mathrm{K}_0(\mathcal{A})$ and apply
	Theorem~\ref{prop:EP} to obtain
	\[ \sum_{p,q\in\Z}(-1)^{p+q}[E_r^{p,q}]
	=\sum_{p,q\in\Z}(-1)^{p+q}[E_{r+1}^{p,q}]. \]

	Lemma~\ref{prop:finite-degeneration} shows that the spectral sequence
	eventually degenerates. Thus, for $r\gg0$,
	$E_r^{p,q}=E_\infty^{p,q}$ for every $(p,q)$. Consequently,
	\[ \sum_{p,q}(-1)^{p+q}[E_r^{p,q}]
	=\sum_{p,q}(-1)^{p+q}[\gr^p\Hm^{p+q}(X)]
	=\sum_n(-1)^n\sum_p[\gr^pH^n], \]
	and the hypotheses ensure that all these sums are finite. By
	Lemma~\ref{prop:K0-prep}~(v) and the properties of
	$\mathrm{F}^\bullet H$, we also have
	$\sum_p[\gr^pH^n(X)]=[\Hm^n(X)]$.
\end{proof}

% segment-id: o014.li2.chapter5.filtered-complex-spectral-sequence.p011
If $(E_r,d_r)_r$ comes from a filtered complex
$(X,d,\mathrm{F}^\bullet)$ and the filtration
$\mathrm{F}^\bullet X^n$ on every $X^n$ is finite, then
Theorem~\ref{prop:classical-convergence} ensures that the strong-convergence
hypothesis $E_r^{p,q}\Rightarrow\Hm^{p+q}(X)$ in
Proposition~\ref{prop:EP-ss} holds automatically. The condition on
$\{(p,q):E_r^{p,q}\neq0\}$, however, is not automatic.

% segment-id: o014.li2.chapter5.filtered-complex-spectral-sequence.p012
Filtered complexes already suffice to produce some simple but useful spectral
sequences; see the exercises for this chapter. Since algebra is the subject of
this book, several of the best-known spectral sequences arise from
bicomplexes. We therefore turn next to the bicomplex case.
